\documentclass[11pt,a4paper]{article}

\usepackage[utf8]{inputenc}
\usepackage[T1]{fontenc}
\usepackage{lmodern}
\usepackage{amsmath, amssymb, amsthm}
\usepackage{mathtools}
\usepackage{geometry}
\usepackage{hyperref}
\usepackage{xcolor}
\usepackage{booktabs}
\usepackage{enumitem}
\usepackage{microtype}
\usepackage{listings}
\usepackage{verbatim}
\usepackage{seqsplit}

\geometry{margin=2.8cm}

\setlength{\emergencystretch}{3em}
\sloppy
\newcommand{\ttbr}[1]{\texttt{\seqsplit{#1}}}
\hypersetup{
  colorlinks=true,
  linkcolor=blue!50!black,
  citecolor=blue!50!black,
  urlcolor=blue!50!black,
  pdftitle={Phantom Orbit Shadowing for the Collatz Conjecture: Verified Core, Two Barriers, and a Repositioning (v4)},
  pdfauthor={Piero Borgatta}
}

\newtheorem{theorem}{Theorem}
\newtheorem{lemma}[theorem]{Lemma}
\newtheorem{proposition}[theorem]{Proposition}
\newtheorem{corollary}[theorem]{Corollary}
\theoremstyle{definition}
\newtheorem{definition}[theorem]{Definition}
\newtheorem{conjecture}[theorem]{Conjecture}
\newtheorem{remark}[theorem]{Remark}
\newtheorem{example}[theorem]{Example}

\DeclareMathOperator{\nutwo}{\nu_2}
\newcommand{\Z}{\mathbb{Z}}
\newcommand{\Zp}{\mathbb{Z}_2}
\newcommand{\Q}{\mathbb{Q}}
\newcommand{\N}{\mathbb{N}}
\newcommand{\R}{\mathbb{R}}
\newcommand{\full}{\mathrm{FULL}}
\newcommand{\core}{\mathrm{CORE}}
\newcommand{\tail}{\mathrm{TAIL}}
\newcommand{\spec}{\rho}

\title{
  Phantom Orbit Shadowing for the Collatz Conjecture: \\
  Verified Core, Two Barriers, and a Repositioning \\
  \large (v4 --- consolidation and limitations)
}
\author{Piero Borgatta\thanks{Independent researcher. Correspondence:
\texttt{info@pieroborgatta.com}; ORCID 0009-0001-8025-2405.}}
\date{\today}

\begin{document}
\maketitle

\begin{abstract}
This is a \emph{consolidation and repositioning} version of an ongoing,
AI-assisted personal research program on the Collatz conjecture. It does
\textbf{not} report progress toward a proof of the conjecture. Its purpose
is the opposite and, we argue, equally worthwhile: to state precisely what
the program has \emph{rigorously established}, to identify exactly
\emph{where} and \emph{why} the program stops, and to record an honest
decision about what to do next.

We first recall the verified core, all mechanically checked in Lean~4 with
Mathlib and \texttt{sorry}-free: an exact congruential shadowing lemma, its
no-infinite-(periodic-)shadowing corollary, a family of finite
Collatz--Wielandt spectral-radius certificates, a finite phantom-set
taxonomy, and an elementary, formally verified bridge reducing the
conjecture to a uniform strict-descent hypothesis.

We then analyze the two analytic branches developed after v3 --- a
weak/operator (averaged) branch and a pointwise first-barrier branch ---
and reach a definite conclusion about each. The weak branch, even if
completed, yields an averaged (``almost-all'') statement and is blocked by
the distributional-to-pointwise barrier. For the pointwise branch we show that, \emph{within the present
phantom/$2$-adic formalism}, the desired
2-adic rigidity on the limit point $\xi_W$ of an aperiodic expanding word
is \emph{logically equivalent} to the nonexistence of divergent integer
orbits, hence cannot act as an independent lever. The reason is a precise
\emph{non-transport lemma} (conditional on the standard Lagarias
parity-vector coding): the Syracuse parity-vector conjugacy carries
eventual periodicity to dynamical pre-periodicity, \emph{not} to
rationality. As a by-product we record one new elementary fact, already
\emph{Lean-checked} in post-prefix congruence form: no positive integer
lies on an expanding cycle (Proposition~\ref{prop:expansive}), generalizing
the no-infinite-shadowing corollary.

Finally we separate the Collatz problem into its two genuinely distinct
hard halves --- absence of divergent orbits (barrier: measure-to-pointwise)
and absence of nontrivial cycles (barrier: Diophantine approximation of
$\log_2 3$) --- locate every part of this program relative to those two
barriers, and state the resulting repositioning: stop pushing the
weak/operator and first-barrier branches as direct routes to the
conjecture; redirect effort to the cycle half (where partial results and
tools genuinely bite) and to Lean formalization, while keeping the
methodology study of AI-assisted mathematics as an explicit goal. All
claims here are honest about scope; we make no claim of subcriticality, no
spectral gap, and no proof of the conjecture.
\end{abstract}

\noindent\textit{Research program note.} This is the fourth versioned
artifact of a personal program by the author, an independent researcher,
exploring whether iterative collaboration with large language model (LLM)
assistants can support non-standard attempts at open problems.
\emph{Testing the practical limits of such collaboration is an explicit
secondary goal, on equal footing with the mathematics.} v4 differs in kind
from v1--v3: where earlier versions added constructions and certificates,
v4 reports the program reaching --- and recognizing --- its barriers, and
documents the resulting strategic decision. We regard the explicit,
provable identification of an approach's own ceiling as a legitimate
scientific output. The earlier versions remain available: concept DOI
\texttt{10.5281/zenodo.20021537}; v2 \texttt{10.5281/zenodo.20098868}; v3
record \url{https://zenodo.org/records/20160154}, DOI
\texttt{10.5281/zenodo.20160154}.

\tableofcontents

\section{Introduction}
\label{sec:intro}

The Collatz conjecture asserts that the map $n \mapsto n/2$ ($n$ even),
$n \mapsto 3n+1$ ($n$ odd) sends every positive integer to $1$. We work
throughout with the accelerated Syracuse map on positive odd integers,
\[
  S(n) = \frac{3n+1}{2^{\nutwo(3n+1)}},
\]
where $\nutwo$ is the $2$-adic valuation; $S$ maps positive odds to
positive odds, and Collatz for $S$ is equivalent to the classical
statement (Section~\ref{sec:bridge}).

\paragraph{What this program built (v1--v3).}
The program grew from a single empirical observation: ``rebel'' orbits of
$S$ shadow, for finite stretches, periodic rational points of the $2$-adic
dynamics attached to negative rational fixed points --- \emph{phantom
cycles} in the terminology of \cite{chang2026}, the underlying object
being classical \cite{lagarias1985}. From this, v1--v3 constructed an
episode graph on phantom orbits, encoded its critical strongly connected
component (SCC) as a finite weighted transfer operator on a refined phase
state, proved an exact congruential shadowing lemma, and produced finite
Collatz--Wielandt spectral-radius certificates --- several of them
mechanically verified in Lean~4. v3 \cite{borgatta_v3} framed the result
as a \emph{conditional spectral reduction}: \emph{if} two explicit
finite-matrix estimates held (a uniform bound under refinement and a
signature-closure condition), the critical SCC would be subcritical.

\paragraph{What changed after v3.}
Substantial post-v3 work pursued those two estimates along two distinct
routes, and this version reports the honest outcome. Neither route reaches
the conjecture, and --- crucially --- we can now say \emph{why}, in each
case precisely. The contributions of v4 are:
\begin{enumerate}[leftmargin=2em]
  \item A clean restatement of the verified core
        (Section~\ref{sec:core}) and of the elementary descent bridge
        (Section~\ref{sec:bridge}), including the observation that the
        bridge \emph{reorganizes} the difficulty rather than reducing it.
  \item A status report on the two analytic branches
        (Section~\ref{sec:branches}): a weak/operator (A0) branch, whose
        finite skeleton is partly formalized but whose natural conclusion
        is an averaged statement; and a pointwise first-barrier branch,
        which targets the correct (every-$n$) statement.
  \item The main new result, a \emph{negative} resolution of the
        currently available phantom/$2$-adic rigidity route
        (Section~\ref{sec:gate}): a non-transport lemma (conditional on
        the standard Lagarias coding) showing that the sought aperiodic
        rigidity is logically equivalent to the divergence half of the
        conjecture itself.
  \item A small new fact (Proposition~\ref{prop:expansive}), Lean-checked
        in post-prefix congruence form: no positive integer lies on an
        expanding cycle.
  \item A precise two-barrier map of the whole problem
        (Section~\ref{sec:barriers}) and the resulting repositioning
        (Section~\ref{sec:decision}).
\end{enumerate}

\paragraph{Stance.}
We make no claim of progress toward Collatz. We claim a verified core, a
provable account of where two natural approaches stop, and a defensible
plan. Throughout, ``rigorous'' means proved (and, where stated, machine
checked); ``empirical'' means numerically supported only; ``conditional''
means valid under a named, unproved hypothesis.

\section{The verified core}
\label{sec:core}

All statements in this section are formalized in Lean~4 with Mathlib and
are \texttt{sorry}-, \texttt{admit}- and \texttt{axiom}-free in the project
sources.

\subsection{Exact congruential shadowing}

For a phantom word $w$ (a nonempty list of positive exponents) with period
sum $A$ and length $L$, let $q_w \in \Q$ be its rational fixed point and
$\xi_m$ its $2$-adic prefix data.

\begin{lemma}[Exact shadowing; \ttbr{exact\_shadowing}]
\label{lem:shadowing}
A positive integer $n$ realizes the first $m$ exponents of $w$ if and only
if $n \equiv q_w \pmod{2^{B_m+1}}$, where $B_m$ is the corresponding
partial period sum. In particular, $b$-fold shadowing is governed by the
single residue condition $n \equiv q_w \pmod{2^{bA+1}}$.
\end{lemma}

This is the iterated step-by-step form of the per-block $2$-adic repulsion
identity of \cite[Prop.~7.4]{chang2026}, equivalent in content to the
classical inverse formula for parity vectors \cite{rozier2025}. We
emphasize this equivalence: the $2$-adic/phantom encoding is a faithful
\emph{re-encoding} of the parity-vector picture, not new information --- a
point that becomes decisive in Section~\ref{sec:gate}.

\begin{corollary}[No infinite periodic shadowing;
\ttbr{no\_infinite\_period\_congruence\_expansive}]
\label{cor:no-infinite}
For an expanding phantom $w$ (i.e.\ $3^L > 2^A$), no positive integer
shadows $w^\infty$ for all $m$: the limiting congruence forces $n = q_w$,
and $q_w < 0$.
\end{corollary}

\subsection{Finite spectral-radius certificates}

For each truncation depth $T$ and high-bit lift count $j$, the program
builds a finite matrix $\full_{T,j} = \core_{T,j} + \tail_{T,j}$ encoding
first-return behavior at the critical node, and applies a weighted
Collatz--Wielandt bound \cite{lemmensnussbaum}. Two certificates are
exposed to Mathlib's real \texttt{spectralRadius}:

\begin{theorem}[\ttbr{t10j32HighBitTailSpectralRadiusBound\_97\_2000}]
$\spec(\full_{10,32}) \le 97/2000 = 0.0485$, with no \texttt{axiom},
\texttt{sorry}, or \texttt{admit} in the generated certificate.
\end{theorem}

\begin{theorem}[\ttbr{k16s16KDeterministicGeneratedSpectralRadiusBound}]
The deterministic finite-residue-cell retention matrix on the compressed
$(K,b)$ macro-state space of the $K_0 = 16$ taxonomy SCC ($37$ states) has
exact rational max ratio
\[
  \frac{90833233962213}{129559208330288} < \frac{3}{4},
\]
with sensitivity confirmed at $\mathrm{lift\_bits} \in \{4,5,6\}$.
\end{theorem}

These are exact statements about \emph{declared finite matrices}. They do
not assert a spectral gap for any infinite operator, and v4 retracts no
part of v3 here: it only stresses the gap between a finite certificate and
an infinite conclusion (Sections~\ref{sec:branches}--\ref{sec:gate}).

\section{The descent bridge, and why it does not reduce difficulty}
\label{sec:bridge}

Let $S^{[k]}$ denote $k$-fold iteration of $S$.

\begin{definition}[\ttbr{UniformStrictDescentHypothesis}]
\label{def:usd}
For every positive odd $n \ne 1$ there exists $k$ with $S^{[k]}(n)$
positive, odd, and $S^{[k]}(n) < n$.
\end{definition}

\begin{theorem}[Lean-verified bridge;
\ttbr{classicalCollatz\_of\_uniformStrictDescent\_provedBridge}]
\label{thm:bridge}
Definition~\ref{def:usd} implies the accelerated Collatz statement (by
strong induction), and the accelerated statement implies the classical
Collatz conjecture (by expanding each $S$-step into one $3n+1$ step
followed by the exact halvings). Both implications are proved in Lean,
\texttt{sorry}-free.
\end{theorem}

\begin{remark}[The bridge reorganizes; it does not reduce]
\label{rem:bridge}
Definition~\ref{def:usd} is the classical ``every $n$ eventually descends
below itself,'' which is equivalent to the convergence half of Collatz.
Theorem~\ref{thm:bridge} is therefore a faithful, mechanically verified
\emph{restatement}, valuable for organizing a formal proof, but it carries
the full difficulty. Every elaboration of the bridge in our sources (a
global descent cover, finite-model soundness interfaces, branch/loss
audits) shares this property: each requires one to \emph{produce} actual
strict-descent witnesses for every source class. Such scaffolding isolates
the obligation cleanly; it does not lighten it. We state this plainly to
forestall the most common misreading of programs of this type.
\end{remark}

\section{The two analytic branches after v3}
\label{sec:branches}

The two estimates left open by v3 were attacked along two routes. We
describe each and its honest status.

\subsection{Branch A: the weak/operator (A0) route}
\label{sec:branchA}

Branch A treats the generated matrices as an averaged row-source kernel
$K_N$ on a finite phase alphabet, with source weight $\mu_N$, and studies
the weighted row-$L^1$ prefix drift
\[
  D_N = \sum_i \mu_N(i) \sum_j \bigl| K_{2N}(i,j) - K_N(i,j) \bigr|.
\]
The target is an A0 weak approximation theorem with named hypotheses
A0W1--A0W7 (finite kernel identification; $D_N \to 0$; an
exceptional high-$\nutwo$ source tail; tail/killing control; uniform
projection and inclusion bounds for a Banach pair $B_s \to B_w$; a limit
operator convention). Genuine, partly formalized, progress was made:

\begin{itemize}[leftmargin=2em]
  \item \textbf{High-$\nutwo$ tail mass, exactly, in Lean}
        (\ttbr{TailCount.dyadic\_tail\_count\_mul\_le},
        \ttbr{dyadic\_tail\_mass\_le}): for the two-prefix source model,
        $\mu_N(\nutwo \ge R) \le 2^{-R}$. This closes one half of the
        $D_N$ split rigorously.
  \item \textbf{A bit-length boundary mechanism, in Lean}
        (\ttbr{biAffineDelta\_dyadicWeight\_period\_boundary\_le} and the
        \texttt{BitLength} development): under a period shift, the
        discrepancy of the dyadic return weights $2^{-\delta}$ is
        controlled by the weighted mass of an explicit set of dyadic
        boundary crossings of cardinality $O(c\log N)$.
  \item \textbf{Exact $2$-adic periodicity of bounded words}
        (script-level, finite): the pure bounded valuation word is
        eventually periodic modulo $2^M$, so the bounded congruential part
        of $D_N$ cancels \emph{exactly} (top dyadic Haar coefficient $0$).
  \item The finite algebraic bridge
        (\ttbr{weighted\_action\_diff\_le}): weighted row-$L^1$ drift
        controls the finite $\ell^\infty \to L^1(\mu_N)$ action error.
\end{itemize}

\begin{remark}[Why Branch A cannot reach Collatz]
\label{rem:branchA}
Every object of Branch A is an \emph{averaged} quantity: $L^1(\mu_N)$,
weighted row means, density-type boundary estimates ($O(\log N / N)$). The
uniform / supremum analogues do not contract (the diagnostics show
$p_{95}$ and maxima persisting; the natural $\Z_2$ martingale Banach space
fails). A completed A0 theorem would establish that the finite matrices
have a legitimate \emph{weak} operator limit; by the very nature of an
averaged norm, the dynamical consequence is an ``almost-all'' / density
statement, not control of every orbit. This is the
distributional-to-pointwise barrier that \cite[Sect.~9]{chang2026}
identifies as the irreducible obstruction to all standard routes. We do
not invoke Hennion or Keller--Liverani: the prerequisites (Banach pair,
boundedness, compactness, a Lasota--Yorke inequality) are absent. Branch A
is therefore retained, if at all, as a \emph{separate} operator-approximation
program, not as a route to the conjecture.
\end{remark}

\subsection{Branch B: the pointwise first-barrier route}
\label{sec:branchB}

Branch B is genuinely different in kind: it is pointwise and parametric,
not averaged. On the source cylinder $n = 103 + 256t$, a parametric
common prefix $[1,1,2,1,1,1]$ (a congruence valid for all $t$, Lean:
\ttbr{a0SemanticDropCommonPrefix\_matches\_source\_cylinder}) carries the
orbit to the endpoint coordinate $593 + 1458t$. One then asks, for every
$t$, whether some matched suffix from the endpoint drops below the source.
The machinery built and verified is the \emph{easy direction} plus its
arithmetic plumbing:
\begin{itemize}[leftmargin=2em]
  \item if a matched suffix crosses the slope barrier
        $1458\cdot 3^{\mathrm{len}} \le 256\cdot 2^{\mathrm{sum}}$ and a
        threshold inequality holds, then $t$ exits
        (\ttbr{a0EndpointSuffixExit\_of\_suffix\_threshold\_contracting});
  \item the threshold is monotone in $t$ (verify at one $t_0$, then
        extend), and matching \emph{forces} a $2$-adic congruence on $t$
        (\ttbr{a0EndpointSuffixCongruence\_of\_matches}).
\end{itemize}
The descent of the whole cylinder then follows --- but only
\emph{conditionally}:

\begin{theorem}[Conditional cylinder descent;
\ttbr{A0EndpointSuffixExitAll\_of\_firstBarrier}]
\label{thm:firstbarrier}
If \ttbr{A0FirstBarrierExistsAll} (every $t$ reaches a first
barrier-crossing suffix) and \ttbr{A0FirstBarrierThresholdAutomatic} (the
threshold holds at that crossing), then every $t$ exits, and every
$n = 103 + 256t$ has a strict descent.
\end{theorem}

Both hypotheses are \texttt{def \ldots\ : Prop} in the sources --- named,
\emph{not} discharged. \ttbr{A0FirstBarrierExistsAll} says that every
orbit on the cylinder eventually contracts past the slope barrier, i.e.\
that no orbit is eternally expanding. This is precisely the correct
(every-$t$) target. The question of v4 is whether the first-barrier
structure supplies a new lever on it.

\section{Negative resolution of the present phantom/2-adic rigidity route}
\label{sec:gate}

We show that the \emph{currently available} phantom/$2$-adic route does
not supply a new lever, and we prove \emph{why}. We flag the rigor status
explicitly: the expanding-cycle fact (Proposition~\ref{prop:expansive}) is
Lean-checked in congruence form, whereas the non-transport lemma
(Lemma~\ref{lem:nontransport}) and the gate resolution
(Proposition~\ref{prop:gate}) are paper-level arguments \emph{conditional
on the standard Lagarias parity-vector coding}, not formalized.

\begin{definition}[Conjugacy and limit point]
Let $\widetilde{S}: \Zp \to \Zp$ be the continuous extension of $S$, and
let $Q:\Zp \to \prod_{k\ge 1}\{a:a\ge 1\}$, $Q(x)_k = \nutwo(3\,
\widetilde{S}^{k-1}(x)+1)$, be the parity-vector conjugacy (a topological
conjugacy of $\widetilde S$ with the shift). For an infinite word $W$ with
all entries $\ge 1$, let $\xi_W = \lim_m \xi_m \in \Zp$ be the limit of the
prefix representatives. If a positive integer $n$ has valuation word $W$,
then $n = \xi_W$.
\end{definition}

\begin{proposition}[Expanding cycles exclude integers]
\label{prop:expansive}
No positive integer lies on a periodic orbit of $S$ whose period word $v$
is expanding ($3^{|v|} > 2^{A_v}$). In particular no positive integer has
an eventually periodic, eventually expanding valuation word.
\end{proposition}
\begin{proof}
A periodic point with word $v$ equals the fixed point
$q_v = C_v/(2^{A_v} - 3^{|v|})$, with $C_v > 0$ and, in the expanding case,
$2^{A_v} - 3^{|v|} < 0$, so $q_v < 0$. Every iterate of a positive odd
integer under $S$ is a positive integer (\ttbr{syracuse\_pos}). A positive
orbit cannot contain $q_v < 0$. The eventual case follows by applying this
to the tail $\widetilde S^{|u|}(n)$ after a finite prefix $u$.
\end{proof}

This is the exact extent of what the phantom/$2$-adic sign argument gives:
it generalizes Corollary~\ref{cor:no-infinite} from the periodic to the
eventually periodic obstruction. Its post-prefix congruence form is
\emph{Lean-checked}, \texttt{sorry}-free, as
\ttbr{no\_positive\_endpoint\_eventually\_periodic\_expansive\_congruence}
in \texttt{NoInfinite.lean}: after any finite prefix, the positive endpoint
of an integer orbit cannot remain, for all numbers of periods, in the
congruence classes of an expanding phantom period. It deliberately
addresses only the eventually periodic case and does \emph{not} rule out an
aperiodic word shadowing different phantom neighborhoods in succession ---
exactly the gap analyzed next.

\begin{lemma}[Non-transport of the rational/periodic dichotomy; standard Lagarias coding]
\label{lem:nontransport}
Under $Q$, eventual periodicity of $W$ corresponds to \emph{pre-periodicity}
of $\xi_W$ under $\widetilde S$ (the orbit enters a cycle), \emph{not} to
rationality of $\xi_W$. Consequently the implication
\[
  \bigl[\, W \text{ not eventually periodic} \ \Longrightarrow\
  \xi_W \notin \Z_{>0} \,\bigr]
\]
is logically equivalent to the nonexistence of divergent positive-integer
$S$-orbits.
\end{lemma}
\begin{proof}
$Q$ is a topological conjugacy, so $W$ is eventually periodic iff
$\xi_W$ is pre-periodic under $\widetilde S$. The contrapositive of the
displayed implication reads
$[\,\xi_W = n \in \Z_{>0} \Rightarrow W \text{ eventually periodic}\,]$,
i.e.\ $[\,n \in \Z_{>0} \Rightarrow n \text{ is pre-periodic}\,]$, i.e.\
every integer orbit eventually enters a cycle --- the negation of the
existence of a divergent integer orbit.
\end{proof}

\begin{remark}[Contrast with binary digits]
For the binary-digit map $d$, ``aperiodic digits $\Rightarrow$ irrational''
is a free theorem, because $d$ transports eventual periodicity to
\emph{rationality}. The valuation/parity vector is a different coordinate
system: $Q$ transports eventual periodicity to \emph{pre-periodicity}, and
``integer $\Rightarrow$ pre-periodic'' is exactly the open divergence
problem. This is the precise structural reason no Diophantine rigidity on
$\xi_W$ is available in the aperiodic case.
\end{remark}

\begin{proposition}[Gate resolution: no new lever within the present formalism]
\label{prop:gate}
On its hard (aperiodic) part, \ttbr{A0FirstBarrierExistsAll} carries the
full classical difficulty of the divergence half of Collatz, with no
additional $2$-adic leverage supplied by the \emph{currently available}
phantom/first-barrier formalism. Concretely:
\begin{enumerate}[leftmargin=2em,label=(\roman*)]
  \item The eventually periodic expanding case is excluded by the sign
        argument (Proposition~\ref{prop:expansive}).
  \item The aperiodic case admits no rationality / irrationality
        obstruction on $\xi_W$ (Lemma~\ref{lem:nontransport}).
  \item Orbit-wide integrality is automatic
        (\ttbr{syracuse\_pos}/\ttbr{syracuse\_odd}) and adds no constraint.
  \item The only genuinely new datum of the formalism --- the sign of the
        periodic fixed point --- is vacuous on both open cases: contractive
        cycles have $q_w > 0$ (no sign contradiction; the cycle problem is
        untouched), and aperiodic words have no fixed-point equation at all
        (the divergence problem is untouched).
\end{enumerate}
\end{proposition}

\begin{remark}[Consistency with the historical record]
This is not surprising; it is the precise articulation, in the program's
own $\xi_W$ language, of why the problem has resisted since the 1970s.
Density-one stopping-time results \cite{terras1976, everett1977} are
reachable because the separating tool is statistical \emph{drift} (the
mean valuation $2 > \log_2 3$); ``every $n$'' is out of reach because drift
does not constrain an individual $\xi_W$. Both branches of this program hit
\emph{this same wall}: Branch A as the distributional-to-pointwise barrier,
Branch B as the absence of rigidity on $\xi_W$ --- two faces of
``density one yes, every $n$ no.''
\end{remark}

\section{The two barriers, precisely}
\label{sec:barriers}

It is useful to separate the conjecture into its two genuinely distinct
hard halves, and to locate every part of this program relative to each.

\paragraph{Half 1: no nontrivial cycles.}
A cycle is a positive-integer solution of $n = C_w/(2^{A} - 3^{L})$; for
$n>0$ one needs $2^A > 3^L$ (contractive). Ruling out all such reduces to
lower bounds on $|A\log 2 - L\log 3|$ for all $(A,L)$ --- the Diophantine
approximation of $\log_2 3$. Baker-type bounds on linear forms in
logarithms give partial exclusions: no $1$-cycles \cite{steiner1977}, no
$2$-cycles \cite{simons2005}, no $m$-cycles up to $m \approx 68$
\cite{simonsdeweger2005}; combined with computational verification to
$\approx 2^{69}$ \cite{barina2021}, any nontrivial cycle must be
astronomically large. The full statement is \emph{open}: the available
bounds leave an uncovered range, and closing it would need a major advance
in the irrationality measure of $\log_2 3$ or a new structural idea. This
is a real but \emph{different} wall from Half~2.

\paragraph{Half 2: no divergent orbits.}
This is the measure-to-pointwise barrier of
Section~\ref{sec:gate}/Remark~\ref{rem:branchA}: density-one descent is
known \cite{terras1976, tao2019}; every-$n$ descent is not, and no
pointwise obstruction is known. This program's phantom/spectral/first-barrier
machinery lives entirely here.

\paragraph{Locating the program.}
\begin{center}
\begin{tabular}{lll}
\toprule
component & half & status \\
\midrule
shadowing lemma + Lean core & both (structure) & rigorous \\
finite CW certificates & --- (finite) & rigorous, finite only \\
expanding-cycle exclusion (Prop.~\ref{prop:expansive}) & Half 1 & rigorous (narrow) \\
weak/operator branch A & Half 2 & blocked (averaged) \\
first-barrier branch B & Half 2 & blocked (Prop.~\ref{prop:gate}) \\
\bottomrule
\end{tabular}
\end{center}

The program contributes a verified \emph{structural} core to both halves
and one narrow rigorous exclusion in Half~1, but its analytic effort is
concentrated entirely in Half~2, against the harder of the two walls.

\section{Decision and planned direction}
\label{sec:decision}

We record the strategic decision reached, with reasons.

\begin{enumerate}[leftmargin=2em]
  \item \textbf{Stop} treating Branch A (weak/operator) and Branch B
        (first-barrier) as routes to the conjecture. By
        Remark~\ref{rem:branchA} and Proposition~\ref{prop:gate}, further
        finite refinement (more $T$, more $j$, more cylinders, more
        suffixes) cannot cross the Half-2 wall \emph{by this route}; it
        only re-confirms the reformulation. Branch A may continue \emph{only} as a separate,
        honestly labeled operator-approximation study.
  \item \textbf{Redirect} mathematical effort to Half~1 (cycles), the only
        half where partial theorems exist and tools bite. The realistic,
        finishable target for this program is the \emph{structural} layer
        (the cycle equation, the Diophantine condition $A/L \approx
        \log_2 3$, elementary exclusions, the link to computational
        verification), \emph{formalized in Lean}. We are explicit about
        the limit: a usable form of Baker's theorem on linear forms in
        logarithms is, to our knowledge, not in Mathlib, so the deep
        analytic exclusions are not cheaply formalizable; the structural
        layer is.
  \item \textbf{Consolidate} the verified assets (this paper) so the
        program's real output --- a \texttt{sorry}-free core, finite
        certificates, a taxonomy, and a precise barrier map --- is on the
        record independently of the conjecture.
  \item \textbf{Continue} the methodology study (Section~\ref{sec:method})
        as an explicit goal: this phase is itself a data point on how far
        cross-AI collaboration carries a non-standard attempt, including
        its capacity to recognize and prove its own ceiling.
\end{enumerate}

\paragraph{What a genuine new idea would have to supply.}
For Half~2, a single-orbit (non-measure) control of the valuation average
--- outside the phantom/$2$-adic circle, since
Lemma~\ref{lem:nontransport} shows that circle is closed. For Half~1, an
improvement in the effective irrationality measure of $\log_2 3$ beyond
current Baker-type bounds, or a structural substitute. We claim neither.

\section{Methodology: AI-assisted attempt on an open problem}
\label{sec:method}

The program is conducted by the author in collaboration with LLM
assistants used in a cross-checking configuration (independent assistants
proposing and auditing each other's constructions, with all formal claims
funneled through Lean). v1--v3 demonstrated that this configuration can
produce nontrivial, machine-verified mathematics (the shadowing core and
the finite certificates) and can manage a large, multi-month formalization.

v4 reports a different and, we think, more informative methodological
outcome: in this phase the cross-AI process \emph{converged on a precise
negative result about its own approach} --- the two-barrier map and the
gate resolution of Section~\ref{sec:gate} --- rather than on new progress.
We regard the ability of such a collaboration to (i) avoid overclaiming,
(ii) localize its own obstruction to a single named hypothesis, and
(iii) prove that hypothesis equivalent to the open problem, as a positive
methodological finding, distinct from progress on the mathematics. The
session logs and the four updated working documents
(\texttt{TODO.md}, \texttt{INVENTORY.md},
\texttt{phase10\_repro\_manifest.md}, \texttt{phase10\_master\_report.md})
record the decision trail.

\section{Honest scope and conclusion}
\label{sec:conclusion}

\paragraph{Established (rigorous, mechanically verified).}
The shadowing lemma and its no-infinite-periodic-shadowing corollary; the
finite Collatz--Wielandt certificates exposed to Mathlib's
\texttt{spectralRadius}; the elementary descent bridge; the verified-core
\texttt{sorry}-freeness; and, new in this version, the post-prefix
congruence form of Proposition~\ref{prop:expansive}
(\ttbr{no\_positive\_endpoint\_eventually\_periodic\_expansive\_congruence}).

\paragraph{Established (on paper, this version; conditional on the standard
Lagarias coding).}
Lemma~\ref{lem:nontransport} (non-transport) and
Proposition~\ref{prop:gate}: within the currently available phantom/$2$-adic
formalism, the first-barrier route supplies no new lever on Half~2. These
are not Lean-formalized and are stated modulo the standard parity-vector
conjugacy.

\paragraph{Empirical (not proved).}
All numerical spectral bounds and decay trends of v1--v3; the finite $D_N$
decay of Branch A.

\paragraph{Missing (and, we argue, behind known barriers).}
A proof of either half of Collatz. Half~2 is blocked by the
measure-to-pointwise barrier, into which both analytic branches of this
program provably fall. Half~1 is open behind the Diophantine wall, with
partial results but no full exclusion.

\paragraph{Honest claim.}
We have built and verified a structural core around the Collatz problem, we
have determined precisely where two natural analytic approaches stop and
proved why, and we have repositioned the program accordingly. We do not
claim a proof of the conjecture, nor progress toward one; we claim a clear,
verified, and honestly bounded map of an attempt.

\section*{Reproducibility}
The Lean~4 development, the enumeration and diagnostic scripts, and the
working notes referenced here are provided in the supplementary archive.
Theorem-to-file correspondences are given inline by Lean declaration name.
Earlier versions: concept DOI \texttt{10.5281/zenodo.20021537}; v3 record
\url{https://zenodo.org/records/20160154}.

\begin{thebibliography}{99}

\bibitem{borgatta_v3}
P.~Borgatta,
\emph{A Spectral Reduction of the Collatz Conjecture via Phantom Orbit
Shadowing} (v3),
Zenodo, DOI \texttt{10.5281/zenodo.20160154}, 2026.

\bibitem{chang2026}
(Chang),
\emph{Phantom cycles and information-theoretic bounds for the Collatz
map}, 2026. (Concurrent, independent work; cited as in v3.)

\bibitem{lagarias1985}
J.~C.~Lagarias,
\emph{The $3x+1$ problem and its generalizations},
Amer. Math. Monthly \textbf{92} (1985), 3--23.

\bibitem{rozier2025}
(Rozier),
\emph{Inverse parity-vector formula for the Syracuse map}, 2025.
(Cited as in v3.)

\bibitem{lemmensnussbaum}
B.~Lemmens and R.~Nussbaum,
\emph{Nonlinear Perron--Frobenius Theory},
Cambridge Tracts in Mathematics \textbf{189}, Cambridge Univ. Press, 2012.

\bibitem{tao2019}
T.~Tao,
\emph{Almost all orbits of the Collatz map attain almost bounded values},
Forum Math. Pi \textbf{10} (2022), e12 (arXiv:1909.03562).

\bibitem{terras1976}
R.~Terras,
\emph{A stopping time problem on the positive integers},
Acta Arith. \textbf{30} (1976), 241--252.

\bibitem{everett1977}
C.~J.~Everett,
\emph{Iteration of the number-theoretic function $f(2n)=n$, $f(2n+1)=3n+2$},
Adv. Math. \textbf{25} (1977), 42--45.

\bibitem{steiner1977}
R.~P.~Steiner,
\emph{A theorem on the Syracuse problem},
Proc. 7th Manitoba Conf. on Numerical Mathematics, 1977, 553--559.

\bibitem{simons2005}
J.~L.~Simons,
\emph{On the nonexistence of $2$-cycles for the $3x+1$ problem},
Math. Comp. \textbf{74} (2005), 1565--1572.

\bibitem{simonsdeweger2005}
J.~L.~Simons and B.~M.~M.~de Weger,
\emph{Theoretical and computational bounds for $m$-cycles of the
$3n+1$ problem},
Acta Arith. \textbf{117} (2005), 51--70.

\bibitem{barina2021}
D.~Barina,
\emph{Convergence verification of the Collatz problem},
J. Supercomput. \textbf{77} (2021), 2681--2688.

\bibitem{kellerliverani}
G.~Keller and C.~Liverani,
\emph{Stability of the spectrum for transfer operators},
Ann. Scuola Norm. Sup. Pisa \textbf{28} (1999), 141--152.

\end{thebibliography}

\end{document}
